\chapter{Solutions to Exercises}
\label{app:solutions}

\section{Chapter 1}

\textbf{1.1.} Write
\[
  \omega_{\mathrm{wind}}=P\,dx+Q\,dy,\qquad
  P=\frac{-y}{x^2+y^2},\quad Q=\frac{x}{x^2+y^2}.
\]
Away from the origin,
\[
  \frac{\partial Q}{\partial x}
  =\frac{y^2-x^2}{(x^2+y^2)^2}
  =\frac{\partial P}{\partial y}.
\]
Thus
\[
  d\omega_{\mathrm{wind}}
  =\left(\frac{\partial Q}{\partial x}
  -\frac{\partial P}{\partial y}\right)dx\wedge dy=0.
\]

\begin{center}
\begin{tikzpicture}[>=Latex, every node/.style={font=\small}]
  \node[smallvisual] (form) at (0,0)
    {$\omega_{\mathrm{wind}}=P\,dx+Q\,dy$};
  \node[smallvisual] (test) at (3.8,0)
    {$d\omega=(Q_x-P_y)\,dx\wedge dy$};
  \node[smallvisual] (zero) at (7.2,0)
    {$Q_x=P_y$\\ \(d\omega_{\mathrm{wind}}=0\)};
  \draw[->, line width=0.5pt] (form) -- (test);
  \draw[->, line width=0.5pt] (test) -- (zero);
\end{tikzpicture}
\end{center}

\textbf{1.2.} With \(x=\cos t\), \(y=\sin t\), one has
\[
  dx=-\sin t\,dt,\qquad dy=\cos t\,dt.
\]
Thus
\[
  -y\,dx+x\,dy=\sin^2t\,dt+\cos^2t\,dt=dt,
\]
and the denominator is \(1\).  The integral from \(0\) to \(2\pi\) is
\(2\pi\).

\begin{center}
\begin{tikzpicture}[>=Latex, scale=1.0, every node/.style={font=\small}]
  \draw[->, black!55, line width=0.45pt] (-1.7,0) -- (1.9,0) node[right] {$x$};
  \draw[->, black!55, line width=0.45pt] (0,-1.6) -- (0,1.85) node[above] {$y$};
  \draw[black, line width=0.65pt] (0,0) circle (1.05);
  \fill[black] (0.74,0.74) circle (1.5pt);
  \draw[->, line width=0.55pt] (0.74,0.74) -- (0.32,1.16)
    node[above left] {$\omega_{\mathrm{wind}}(\dot\gamma)=1$};
  \draw[->, line width=0.45pt, black!60] (0,0) -- (0.74,0.74)
    node[midway, below right] {$r$};
  \draw[->, line width=0.45pt] (0.45,0)
    arc[start angle=0,end angle=45,radius=0.45];
  \node at (0.64,0.2) {$\theta$};
  \node[align=center] at (0,-2.05)
    {On \(S^1\), the winding form pulls back to \(d\theta=dt\).};
\end{tikzpicture}
\end{center}

\textbf{1.3.} Write \(f=(z-p)^m u(z)\), where \(u(p)\ne0\).  Then
\[
  \frac{df}{f}=m\,\frac{dz}{z-p}+\frac{du}{u}.
\]
The second term is holomorphic near \(p\), so the residue is \(m=\ord_p(f)\).

\begin{center}
\begin{tikzpicture}[>=Latex, every node/.style={font=\small}]
  \node[smallvisual] (factor) at (0,0)
    {$f=(z-p)^m u$\\ \(u(p)\ne0\)};
  \node[smallvisual] (log) at (3.45,0)
    {$df/f=m\,dz/(z-p)+du/u$};
  \node[smallvisual] (res) at (6.85,0)
    {$\Res_p(df/f)=m$\\ \(=\ord_p(f)\)};
  \draw[->, line width=0.5pt] (factor) -- (log);
  \draw[->, line width=0.5pt] (log) -- (res);
\end{tikzpicture}
\end{center}

\textbf{1.4.} By Exercise 1.3, the residue of \(df/f\) at \(p\) is
\(\ord_p(f)\).  On a compact Riemann surface the sum of residues of a
meromorphic one-form is zero.  Therefore
\[
  0=\sum_p\Res_p(df/f)=\sum_p\ord_p(f)=\degdiv\divisor(f).
\]

\begin{center}
\begin{tikzpicture}[>=Latex, every node/.style={font=\small}]
  \draw[line width=0.6pt] (0,0) ellipse (2.05 and 0.9);
  \foreach \x/\y/\m in {-1.1/0.25/{m_1},0.15/-0.2/{m_2},1.15/0.28/{m_3}}{
    \fill[black] (\x,\y) circle (1.5pt);
    \node[above] at (\x,\y) {$\m$};
  }
  \node[align=center] at (4.6,0.15)
    {$\displaystyle\sum_p\Res_p(df/f)=0$\\[3pt]
     \(\displaystyle\sum_p\ord_p(f)=0\)};
  \draw[->, line width=0.55pt] (2.25,0) -- (3.25,0);
  \node at (0,-1.25) {compact surface};
\end{tikzpicture}
\end{center}

\section{Chapter 2}

\textbf{2.1.} If \(\chi=0\), then \(0=2-2g\), so \(g=1\).

\begin{center}
\begin{tikzpicture}[>=Latex, every node/.style={font=\small}]
  \node[smallvisual] (chi) at (0,0) {$\chi=0$};
  \node[smallvisual] (eq) at (2.4,0) {$0=2-2g$};
  \node[smallvisual] (genus) at (4.8,0) {$g=1$};
  \draw[->, line width=0.5pt] (chi) -- (eq);
  \draw[->, line width=0.5pt] (eq) -- (genus);
\end{tikzpicture}
\end{center}

\textbf{2.2.} Noncompact surfaces need not give compact images, so a maximum of
\(|f|\) may not exist.  For example, \(f(z)=z\) on \(\C\) is holomorphic and
nonconstant.

\begin{center}
\begin{tikzpicture}[>=Latex, every node/.style={font=\small}]
  \draw[->, black!55, line width=0.45pt] (-2.2,0) -- (2.2,0) node[right] {$\C$};
  \draw[->, black!55, line width=0.45pt] (0,-0.7) -- (0,0.85);
  \draw[->, line width=0.55pt] (-1.5,0) -- (1.55,0)
    node[midway, above] {$f(z)=z$};
  \node[align=center] at (0,-1.15)
    {noncompact domain: \(|f|\) can grow without a maximum};
\end{tikzpicture}
\end{center}

\textbf{2.3.} For a nonzero meromorphic function, zeros and poles are isolated.
An infinite subset of a compact surface has an accumulation point.  That would
contradict isolation: a holomorphic function with accumulating zeros is
identically zero locally, and poles are isolated by the local form
\((z-p)^{-m}u(z)\).  Hence the zero and pole sets are finite.

\begin{center}
\begin{tikzpicture}[>=Latex, every node/.style={font=\small}]
  \draw[line width=0.6pt] (0,0) ellipse (2.0 and 0.9);
  \foreach \x/\y in {-1.2/0.25,-0.45/-0.18,0.65/0.22,1.25/-0.2}{
    \fill[black] (\x,\y) circle (1.35pt);}
  \draw[dashed, line width=0.45pt] (1.25,-0.2) circle (0.32);
  \node[align=center] at (4.35,0.15)
    {isolated zeros/poles\\on a compact surface\\must be finite};
  \draw[->, line width=0.55pt] (2.25,0) -- (3.05,0);
\end{tikzpicture}
\end{center}

\textbf{2.4.} A parallelogram with opposite sides identified has one face, two
edge classes, and one vertex class.  Hence
\[
  \chi=1-2+1=0.
\]
Since \(\chi=2-2g\), the quotient has \(g=1\).  Thus the quotient surface is a
torus.

\begin{center}
\begin{tikzpicture}[>=Latex, every node/.style={font=\small}]
  \coordinate (A) at (0,0);
  \coordinate (B) at (2.7,0);
  \coordinate (C) at (3.55,1.45);
  \coordinate (D) at (0.85,1.45);
  \draw[line width=0.6pt] (A)--(B)--(C)--(D)--cycle;
  \draw[->, line width=0.55pt] ($(A)!0.36!(B)$) -- ($(A)!0.64!(B)$);
  \draw[->, line width=0.55pt] ($(D)!0.64!(C)$) -- ($(D)!0.36!(C)$);
  \draw[->, line width=0.55pt] ($(A)!0.35!(D)$) -- ($(A)!0.65!(D)$);
  \draw[->, line width=0.55pt] ($(B)!0.65!(C)$) -- ($(B)!0.35!(C)$);
  \node[below] at (1.35,0) {$\omega_1$};
  \node[left] at (0.42,0.72) {$\omega_2$};
  \draw[->, line width=0.55pt] (4.1,0.72) -- node[above] {identify sides} (5.3,0.72);
  \begin{scope}[xshift=6.9cm, yshift=0.72cm, scale=0.23]
    \drawprecisetorus
  \end{scope}
  \node[align=center] at (6.9,-1.0)
    {one face, two edge classes,\\one vertex class};
\end{tikzpicture}
\end{center}

\section{Chapter 3}

\textbf{3.1.} If \(u'=u+\lambda\) and \(v'=v+\mu\), then
\[
  u'+v'=u+v+\lambda+\mu.
\]
Since \(\lambda+\mu\in\Lambda\), this represents the same class as \(u+v\).

\begin{center}
\begin{tikzpicture}[>=Latex, every node/.style={font=\small}]
  \node[smallvisual] (reps) at (0,0)
    {$u'=u+\lambda$\\ \(v'=v+\mu\)};
  \node[smallvisual] (sum) at (3.6,0)
    {$u'+v'=u+v+\lambda+\mu$};
  \node[smallvisual] (same) at (7.1,0)
    {$[u'+v']=[u+v]$};
  \draw[->, line width=0.5pt] (reps) -- (sum);
  \draw[->, line width=0.5pt] (sum) -- (same);
\end{tikzpicture}
\end{center}

\textbf{3.2.} A matrix in \(\GL_2(\Z)\) gives integer combinations of
\(\omega_1,\omega_2\), and its inverse also has integer entries.  Thus each old
basis vector is an integer combination of the new ones and conversely, so the
generated subgroup is unchanged.

\begin{center}
\begin{tikzpicture}[>=Latex, every node/.style={font=\small}]
  \draw[->, black!55, line width=0.45pt] (0,0) -- (1.8,0) node[below] {$\omega_1$};
  \draw[->, black!55, line width=0.45pt] (0,0) -- (0.65,1.35) node[left] {$\omega_2$};
  \draw[->, line width=0.55pt] (2.4,0.6) -- node[above] {$A\in\GL_2(\Z)$} (4.0,0.6);
  \begin{scope}[xshift=4.6cm]
    \draw[->, black!55, line width=0.45pt] (0,0) -- (1.45,0.55) node[below right] {$\omega'_1$};
    \draw[->, black!55, line width=0.45pt] (0,0) -- (-0.25,1.45) node[left] {$\omega'_2$};
  \end{scope}
  \node[align=center] at (2.95,-0.75)
    {integer change of basis preserves the lattice};
\end{tikzpicture}
\end{center}

\textbf{3.3.} The ratio \(\tau=\omega_2/\omega_1\) is non-real because, if
\(\tau\in\R\), then \(\omega_2=\tau\omega_1\) would make
\(\omega_1,\omega_2\) real-linearly dependent.  If \(\tau=a+ib\) with
\(b<0\), swapping the basis vectors gives the new ratio
\[
  \frac{\omega_1}{\omega_2}=\frac1{\tau}
  =\frac{a-ib}{a^2+b^2},
\]
whose imaginary part is \(-b/(a^2+b^2)>0\).  Hence an oriented basis can be
chosen with \(\tau\in\HH\).

\begin{center}
\begin{tikzpicture}[>=Latex, every node/.style={font=\small}]
  \draw[->, black!55, line width=0.45pt] (-2.2,0) -- (2.2,0) node[right] {$\operatorname{Re}\tau$};
  \draw[->, black!55, line width=0.45pt] (0,-1.25) -- (0,1.45) node[above] {$\operatorname{Im}\tau$};
  \fill[black] (1.0,-0.7) circle (1.5pt) node[below right] {$\tau$};
  \fill[black] (1.0,0.7) circle (1.5pt) node[above right] {$\tau'\in\HH$};
  \draw[->, line width=0.5pt] (1.0,-0.55) -- (1.0,0.55);
  \node[align=center] at (-1.15,0.75)
    {choose orientation\\with positive imaginary part};
\end{tikzpicture}
\end{center}

\textbf{3.4.} The map is well-defined because if \(z'=z+\lambda\), then
\[
  z'+a=z+a+\lambda,
\]
so \([z'+a]=[z+a]\).  It is holomorphic because it is induced locally by the
holomorphic map \(z\mapsto z+a\) on \(\C\).  Its inverse is translation by
\(-a\), so it is biholomorphic.

\begin{center}
\begin{tikzpicture}[>=Latex, every node/.style={font=\small}]
  \coordinate (A) at (0,0);
  \coordinate (B) at (3.0,0);
  \coordinate (C) at (3.75,1.6);
  \coordinate (D) at (0.75,1.6);
  \draw[line width=0.55pt] (A)--(B)--(C)--(D)--cycle;
  \draw[->, line width=0.45pt, black!60] (0.75,0.45) -- (1.7,0.9)
    node[midway, below] {$+a$};
  \fill[black] (0.75,0.45) circle (1.5pt) node[below left] {$z$};
  \fill[black] (1.7,0.9) circle (1.5pt) node[above right] {$z+a$};
  \draw[->, line width=0.45pt, black!60] (0.75,0.45) -- (3.75,0.45)
    node[midway, below] {$+\lambda$};
  \fill[black] (3.75,0.45) circle (1.3pt) node[right] {$z+\lambda$};
  \draw[->, line width=0.45pt, black!60] (3.75,0.45) -- (4.7,0.9)
    node[midway, below right] {$+a$};
  \node[align=center] at (2,-0.9)
    {\([z]=[z+\lambda]\) and \([z+a]=[z+\lambda+a]\).};
\end{tikzpicture}
\end{center}

\textbf{3.5.} For a map \(F:G\to H\), the checklist is:
\[
\begin{array}{c|c}
\text{condition} & \text{calculation or property}\\ \hline
\text{group homomorphism} & F(x+y)=F(x)+F(y)\\
\text{holomorphic map} & F \text{ is holomorphic in local complex charts}\\
\text{holomorphic inverse} & F^{-1} \text{ exists and is holomorphic.}
\end{array}
\]
Thus an isomorphism of complex Lie groups is simultaneously a biholomorphism and
a group isomorphism, not merely an abstract abelian-group isomorphism.

\begin{center}
\begin{tikzpicture}[>=Latex, every node/.style={font=\small}]
  \node[smallvisual] (iso) at (0,0) {complex Lie\\group isomorphism};
  \node[smallvisual] (hol) at (3.0,0.75) {biholomorphic};
  \node[smallvisual] (grp) at (3.0,-0.75) {respects addition};
  \draw[->, line width=0.5pt] (iso) -- (hol);
  \draw[->, line width=0.5pt] (iso) -- (grp);
  \node[align=center] at (5.6,0)
    {stronger than an abstract\\abelian-group isomorphism};
\end{tikzpicture}
\end{center}

\section{Chapter 4}

\textbf{4.1.} A nonconstant elliptic function cannot be holomorphic, so it has a
pole.  If the total pole order were \(1\), there would be exactly one simple
pole.  Its Laurent expansion would have a nonzero coefficient of
\((z-a)^{-1}\), hence a nonzero residue.  This contradicts the residue-sum law
for elliptic functions.  Therefore the total pole order is at least \(2\).

\begin{center}
\begin{tikzpicture}[>=Latex, every node/.style={font=\small}]
  \draw[line width=0.6pt] (0,0) ellipse (1.7 and 0.8);
  \fill[black] (0.25,0.1) circle (1.5pt) node[above] {simple pole};
  \node[smallvisual] (res) at (4.0,0.25) {nonzero residue};
  \node[smallvisual] (bad) at (6.7,0.25) {contradicts\\residue sum \(0\)};
  \draw[->, line width=0.55pt] (1.9,0.1) -- (res);
  \draw[->, line width=0.55pt] (res) -- (bad);
  \node at (0,-1.1) {total pole order \(1\) is impossible};
\end{tikzpicture}
\end{center}

\textbf{4.2.} Choose a period parallelogram whose boundary contains no zeros or
poles of \(f\).  Apply the argument principle to \(f\) on its boundary:
\[
  \frac{1}{2\pi i}\int_{\partial P}\frac{f'}{f}\,dz.
\]
Opposite sides cancel by periodicity.  The integral is therefore zero, and it
equals the number of zeros minus the number of poles, counted with
multiplicity.

\begin{center}
\begin{tikzpicture}[>=Latex, every node/.style={font=\small}]
  \coordinate (A) at (0,0);
  \coordinate (B) at (2.7,0);
  \coordinate (C) at (3.45,1.45);
  \coordinate (D) at (0.75,1.45);
  \draw[line width=0.55pt] (A)--(B)--(C)--(D)--cycle;
  \draw[->, line width=0.5pt] (0.55,0) -- (1.65,0);
  \draw[->, line width=0.5pt] (3.12,1.45) -- (2.05,1.45);
  \draw[->, line width=0.5pt] (0.18,0.35) -- (0.52,1.05);
  \draw[->, line width=0.5pt] (3.28,1.05) -- (2.95,0.35);
  \node[align=center] at (5.6,0.75)
    {periodicity pairs\\opposite boundary integrals\\with opposite signs};
  \draw[->, line width=0.55pt] (3.8,0.75) -- (4.45,0.75);
\end{tikzpicture}
\end{center}

\textbf{4.3.} Changing a lift \(z_i\) by a lattice element changes the weighted
sum by an element of \(\Lambda\).  Therefore the Abel sum is intrinsic only
modulo \(\Lambda\), that is, as a point of \(\C/\Lambda\).

\begin{center}
\begin{tikzpicture}[>=Latex, every node/.style={font=\small}]
  \node[smallvisual] (lift) at (0,0)
    {$z_i\rightsquigarrow z_i+\lambda_i$};
  \node[smallvisual] (sum) at (3.4,0)
    {$\sum m_i z_i\rightsquigarrow
      \sum m_i z_i+\lambda$};
  \node[smallvisual] (quot) at (6.85,0)
    {well-defined only\\in \(\C/\Lambda\)};
  \draw[->, line width=0.5pt] (lift) -- (sum);
  \draw[->, line width=0.5pt] (sum) -- (quot);
\end{tikzpicture}
\end{center}

\textbf{4.4.} If \(F\) is meromorphic on \(T\), then \(F\circ q\) is meromorphic
on \(\C\) and satisfies
\[
  F(q(z+\lambda))=F(q(z))
\]
for all \(\lambda\in\Lambda\).  Conversely, if \(f\) is elliptic, define
\(F([z])=f(z)\).  Periodicity makes this well-defined, and the quotient charts
show that \(F\) is meromorphic on \(T\).

\begin{center}
\begin{tikzpicture}[>=Latex, every node/.style={font=\small}]
  \node[smallvisual] (C) at (0,0.75) {$\C$};
  \node[smallvisual] (T) at (0,-0.75) {$T=\C/\Lambda$};
  \node[smallvisual] (vals) at (4.1,-0.75) {$\CP^1$};
  \draw[->, line width=0.55pt] (C) -- node[left] {$q$} (T);
  \draw[->, line width=0.55pt] (T) -- node[below] {$F$} (vals);
  \draw[->, line width=0.55pt] (C) -- node[above] {$f=F\circ q$} (vals);
  \node[align=center] at (4.1,0.8)
    {periodic meromorphic functions\\are functions on the quotient};
\end{tikzpicture}
\end{center}

\textbf{4.5.} For
\[
  \divisor(f)=[a_1]+[a_2]-[b_1]-[b_2],
\]
the third Liouville law gives
\[
  a_1+a_2-b_1-b_2\in\Lambda.
\]
Modulo \(\Lambda\), this is
\[
  a_1+a_2-b_1-b_2\equiv0\pmod{\Lambda}.
\]
By the definition of the Abel-sum map, this says
\[
  \smap(\divisor(f))=0\in T.
\]

\begin{center}
\begin{tikzpicture}[>=Latex, every node/.style={font=\small}]
  \coordinate (A) at (0,0);
  \coordinate (B) at (3.2,0);
  \coordinate (C) at (4.1,1.7);
  \coordinate (D) at (0.9,1.7);
  \draw[line width=0.55pt] (A)--(B)--(C)--(D)--cycle;
  \draw[->, line width=0.45pt] ($(A)!0.25!(B)$) -- ($(A)!0.55!(B)$);
  \draw[->, line width=0.45pt] ($(D)!0.55!(C)$) -- ($(D)!0.25!(C)$);
  \draw[->, line width=0.45pt] ($(A)!0.3!(D)$) -- ($(A)!0.62!(D)$);
  \draw[->, line width=0.45pt] ($(B)!0.62!(C)$) -- ($(B)!0.3!(C)$);
  \fill[black] (1.2,0.55) circle (1.5pt) node[below] {$a_1,+$};
  \fill[white] (2.7,1.1) circle (1.6pt);
  \draw[black] (2.7,1.1) circle (1.6pt) node[above right] {$b_1,-$};
  \fill[black] (3.25,0.55) circle (1.5pt) node[below right] {$a_2,+$};
  \fill[white] (1.85,1.25) circle (1.6pt);
  \draw[black] (1.85,1.25) circle (1.6pt) node[above left] {$b_2,-$};
  \node[align=center] at (6.2,0.95)
    {$\divisor(f)=[a_1]+[a_2]-[b_1]-[b_2]$\\[2pt]
     \(a_1+a_2-b_1-b_2\equiv0\pmod{\Lambda}\)};
  \draw[->, line width=0.55pt] (4.35,0.9) -- (5.15,0.9);
\end{tikzpicture}
\end{center}

\section{Chapter 5}

\textbf{5.1.} Replacing \(z\) by \(-z\) permutes the lattice terms
\(\lambda\mapsto-\lambda\) and leaves the series unchanged.  Hence
\(\wp(-z)=\wp(z)\).

\begin{center}
\begin{tikzpicture}[>=Latex, every node/.style={font=\small}]
  \draw[line width=0.5pt] (-1.1,0) -- (1.1,0);
  \draw[line width=0.5pt] (0,-0.75) -- (0,0.75);
  \foreach \x in {-1,0,1}{\foreach \y in {-0.6,0,0.6}{\fill[black] (\x,\y) circle (1pt);}}
  \fill[black] (0.55,0.32) circle (1.4pt) node[above right] {$z$};
  \fill[black] (-0.55,-0.32) circle (1.4pt) node[below left] {$-z$};
  \draw[->, line width=0.5pt] (0.4,0.25) -- (-0.4,-0.25);
  \node[align=center] at (3.65,0)
    {the lattice is unchanged by\\\(\lambda\mapsto-\lambda\), so \(\wp\) is even};
\end{tikzpicture}
\end{center}

\textbf{5.2.} The derivative of an even function is odd: differentiating
\(\wp(-z)=\wp(z)\) gives \(-\wp'(-z)=\wp'(z)\), so
\(\wp'(-z)=-\wp'(z)\).

\begin{center}
\begin{tikzpicture}[>=Latex, every node/.style={font=\small}]
  \node[smallvisual] (even) at (0,0) {$\wp(-z)=\wp(z)$};
  \node[smallvisual] (diff) at (3.0,0) {$-\wp'(-z)=\wp'(z)$};
  \node[smallvisual] (odd) at (6.0,0) {$\wp'(-z)=-\wp'(z)$};
  \draw[->, line width=0.5pt] (even) -- node[above] {differentiate} (diff);
  \draw[->, line width=0.5pt] (diff) -- (odd);
\end{tikzpicture}
\end{center}

\textbf{5.3.} The equation
\[
  \wp'^2=4\wp^3-g_2\wp-g_3
\]
is a polynomial equation for \(\wp'\) with coefficients in \(\C(\wp)\).  It is
quadratic in \(\wp'\), so \(\wp'\) has degree at most \(2\) over \(\C(\wp)\).

\begin{center}
\begin{tikzpicture}[>=Latex, every node/.style={font=\small}]
  \node[smallvisual] (field) at (0,0) {$\C(\wp)$};
  \node[smallvisual] (quad) at (3.2,0)
    {$Y^2-(4\wp^3-g_2\wp-g_3)=0$};
  \node[smallvisual] (deg) at (6.6,0)
    {degree of \(\wp'\)\\is at most \(2\)};
  \draw[->, line width=0.5pt] (field) -- (quad);
  \draw[->, line width=0.5pt] (quad) -- (deg);
\end{tikzpicture}
\end{center}

\textbf{5.4.} From the Laurent expansion
\[
  \wp(z)=z^{-2}+O(z^2),
\]
\(\wp\) has a pole of order \(2\) at \(O\).  From
\[
  \wp'(z)=-2z^{-3}+O(z),
\]
\(\wp'\) has a pole of order \(3\) at \(O\).

\begin{center}
\begin{tikzpicture}[>=Latex, every node/.style={font=\small}]
  \node[smallvisual] (wp) at (0,0.55)
    {$\wp=z^{-2}+O(z^2)$\\ pole order \(2\)};
  \node[smallvisual] (wpp) at (4.1,0.55)
    {$\wp'=-2z^{-3}+O(z)$\\ pole order \(3\)};
  \node[align=center] at (2.05,-0.75)
    {the leading negative power gives the pole order at \(O\)};
  \draw[->, line width=0.5pt] (wp) -- (wpp);
\end{tikzpicture}
\end{center}

\textbf{5.5.} Substituting \(x=\wp(z)\) and \(y=\wp'(z)\) into the differential
equation gives
\[
  y^2=4x^3-g_2x-g_3.
\]
Thus every point \((\wp(z),\wp'(z))\), away from the pole, lies on the affine
cubic.

\begin{center}
\begin{tikzpicture}[>=Latex, every node/.style={font=\small}]
  \node[smallvisual] (torus) at (0,0) {$z\in\C/\Lambda$};
  \node[smallvisual] (coord) at (3.0,0) {$(x,y)=(\wp(z),\wp'(z))$};
  \node[smallvisual] (curve) at (6.6,0) {$y^2=4x^3-g_2x-g_3$};
  \draw[->, line width=0.5pt] (torus) -- (coord);
  \draw[->, line width=0.5pt] (coord) -- (curve);
\end{tikzpicture}
\end{center}

\section{Chapter 6}

\textbf{6.1.} The function \(z^2(z-1)^{-3}\) has a zero of order \(2\) at
\(0\), a pole of order \(3\) at \(1\), and a pole/zero contribution at
\(\infty\) chosen so the total degree is zero.  In the coordinate \(w=1/z\),
the function is
\[
  \frac{w}{(1-w)^3},
\]
so it has a zero of order \(1\) at \(\infty\).  Hence
\[
  \divisor(z^2(z-1)^{-3})=2[0]-3[1]+[\infty].
\]

\begin{center}
\begin{tikzpicture}[>=Latex, every node/.style={font=\small}]
  \draw[->, black!55, line width=0.45pt] (-0.5,0) -- (6.4,0);
  \foreach \x/\lab/\mult in {1/{0}/{+2},3/{1}/{-3},5.4/{\infty}/{+1}}{
    \fill[black] (\x,0) circle (1.5pt);
    \node[above] at (\x,0) {$\lab$};
    \node[below] at (\x,0) {$\mult$};
  }
  \node[align=center] at (3,-1.05)
    {zeros and poles balance to total degree \(0\)};
\end{tikzpicture}
\end{center}

\textbf{6.2.} Products of meromorphic functions add divisors:
\[
  \divisor(fg)=\divisor(f)+\divisor(g).
\]
Inverses negate divisors, and constants give the zero divisor.  Since principal
divisors have degree zero, \(\Prin(X)\) is a subgroup of \(\Div^0(X)\).

\begin{center}
\begin{tikzpicture}[>=Latex, every node/.style={font=\small}]
  \node[smallvisual] (f) at (0,0.6) {$f$};
  \node[smallvisual] (g) at (0,-0.6) {$g$};
  \node[smallvisual] (fg) at (3.0,0) {$fg$};
  \node[smallvisual] (div) at (6.25,0)
    {$\divisor(fg)=\divisor(f)+\divisor(g)$};
  \draw[->, line width=0.5pt] (f) -- (fg);
  \draw[->, line width=0.5pt] (g) -- (fg);
  \draw[->, line width=0.5pt] (fg) -- (div);
\end{tikzpicture}
\end{center}

\textbf{6.3.} The coefficients of \([P]-[O]\) are \(1\) and \(-1\), so the
degree is \(1-1=0\).

\begin{center}
\begin{tikzpicture}[>=Latex, every node/.style={font=\small}]
  \node[smallvisual] (P) at (0,0) {$[P]$\\ \(+1\)};
  \node[smallvisual] (O) at (2.3,0) {$[O]$\\ \(-1\)};
  \node[smallvisual] (zero) at (4.7,0) {$1-1=0$\\ degree zero};
  \draw[->, line width=0.5pt] (P) -- (O);
  \draw[->, line width=0.5pt] (O) -- (zero);
\end{tikzpicture}
\end{center}

\textbf{6.4.} The divisor
\[
  [P]+[Q]-2[O]
\]
has degree \(1+1-2=0\).  Therefore it lies in \(\Div^0(X)\), and its class
modulo principal divisors is an element of \(\Pic^0(X)\).

\begin{center}
\begin{tikzpicture}[>=Latex, every node/.style={font=\small}]
  \node[smallvisual] (div) at (0,0)
    {$[P]+[Q]-2[O]$\\ degree \(0\)};
  \node[smallvisual] (Div0) at (3.2,0) {$\Div^0(X)$};
  \node[smallvisual] (Pic) at (6.1,0) {$\Pic^0(X)$};
  \draw[->, line width=0.5pt] (div) -- (Div0);
  \draw[->, line width=0.5pt] (Div0) -- node[above] {mod \(\Prin(X)\)} (Pic);
\end{tikzpicture}
\end{center}

\textbf{6.5.} Two degree-zero divisors \(D,D'\) define the same element of
\(\Pic^0(X)=\Div^0(X)/\Prin(X)\) exactly when
\[
  D-D'\in\Prin(X).
\]
This is the defining equivalence relation in the quotient group.

\begin{center}
\begin{tikzpicture}[>=Latex, every node/.style={font=\small}]
  \node[smallvisual] (D) at (0,0.65) {$D$};
  \node[smallvisual] (Dp) at (0,-0.65) {$D'$};
  \node[smallvisual] (diff) at (3.0,0) {$D-D'\in\Prin(X)$};
  \node[smallvisual] (same) at (6.0,0) {$[D]=[D']$\\in \(\Pic^0(X)\)};
  \draw[->, line width=0.5pt] (D) -- (diff);
  \draw[->, line width=0.5pt] (Dp) -- (diff);
  \draw[->, line width=0.5pt] (diff) -- (same);
\end{tikzpicture}
\end{center}

\section{Chapter 7}

\textbf{7.1.} The dimensions are
\[
  \ell(O)=1,\qquad \ell(2O)=2,\qquad \ell(3O)=3.
\]

\begin{center}
\begin{tikzpicture}[>=Latex, every node/.style={font=\small}]
  \node[smallvisual] (one) at (0,0) {$O$\\ \(\ell=1\)};
  \node[smallvisual] (two) at (2.7,0) {$2O$\\ \(\ell=2\)};
  \node[smallvisual] (three) at (5.4,0) {$3O$\\ \(\ell=3\)};
  \draw[->, line width=0.5pt] (one) -- (two);
  \draw[->, line width=0.5pt] (two) -- (three);
  \node[align=center] at (2.7,-1.0)
    {genus-one Riemann--Roch: \(\ell(nO)=n\)};
\end{tikzpicture}
\end{center}

\textbf{7.2.} A function in \(L(O)\) has at most a simple pole at \(O\) and no
other poles.  If it has no pole, it is holomorphic on the compact torus and is
constant.  If it has a simple pole at \(O\), it has total pole order \(1\),
which is impossible by Exercise 4.1.  Hence every element of \(L(O)\) is
constant.

\begin{center}
\begin{tikzpicture}[>=Latex, every node/.style={font=\small}]
  \node[smallvisual] (LO) at (0,0) {$f\in L(O)$};
  \node[smallvisual] (simple) at (3.0,0.65) {simple pole?\\total order \(1\)};
  \node[smallvisual] (no) at (6.0,0.65) {impossible};
  \node[smallvisual] (const) at (3.0,-0.75) {no pole\\compact holomorphic};
  \node[smallvisual] (C) at (6.0,-0.75) {constant};
  \draw[->, line width=0.5pt] (LO) -- (simple);
  \draw[->, line width=0.5pt] (simple) -- (no);
  \draw[->, line width=0.5pt] (LO) -- (const);
  \draw[->, line width=0.5pt] (const) -- (C);
\end{tikzpicture}
\end{center}

\textbf{7.3.} The pole orders at \(O\) are \(0,2,3\) for \(1,\wp,\wp'\).  In a
linear relation, the term with the largest pole order cannot be cancelled by
terms with smaller pole order.  Descending through the possible largest pole
orders forces all coefficients to be zero.

\begin{center}
\begin{tikzpicture}[>=Latex, every node/.style={font=\small}]
  \draw[->, black!55, line width=0.45pt] (-0.3,0) -- (5.2,0) node[right] {pole order};
  \foreach \x/\lab in {0/{1},2/{\wp},3/{\wp'}}{
    \fill[black] (\x,0) circle (1.4pt);
    \node[above] at (\x,0) {$\lab$};
    \node[below] at (\x,0) {$\x$};
  }
  \node[align=center] at (2.6,-1.1)
    {the largest pole order in a relation cannot be cancelled};
\end{tikzpicture}
\end{center}

\textbf{7.4.} The constant function satisfies
\[
  \divisor(1)+2[O]=2[O]\ge0,
\]
so \(1\in L(2O)\).  Since \(\wp\) has no poles away from \(O\) and has pole
order \(2\) at \(O\),
\[
  \divisor(\wp)+2[O]\ge0,
\]
so \(\wp\in L(2O)\).  The pole orders \(0\) and \(2\) make \(1,\wp\) linearly
independent.  Since \(\ell(2O)=2\), they form a basis:
\[
  L(2O)=\Span_\C(1,\wp).
\]

\begin{center}
\begin{tikzpicture}[>=Latex, every node/.style={font=\small}]
  \node[smallvisual] (ind) at (0,0) {$1,\wp$\\independent};
  \node[smallvisual] (dim) at (3.0,0) {$\ell(2O)=2$};
  \node[smallvisual] (basis) at (6.0,0) {$L(2O)=\Span_\C(1,\wp)$};
  \draw[->, line width=0.5pt] (ind) -- (dim);
  \draw[->, line width=0.5pt] (dim) -- (basis);
\end{tikzpicture}
\end{center}

\textbf{7.5.} Membership is checked by the divisor inequalities
\[
  \divisor(1)+3[O]=3[O]\ge0,\qquad
  \divisor(\wp)+3[O]\ge0,\qquad
  \divisor(\wp')+3[O]\ge0.
\]
The second inequality holds because \(\wp\) has pole order \(2\le3\) at \(O\);
the third holds because \(\wp'\) has pole order exactly \(3\) at \(O\).  The
functions \(1,\wp,\wp'\) are linearly independent by pole orders \(0,2,3\).
Since \(\ell(3O)=3\), they form a basis:
\[
  L(3O)=\Span_\C(1,\wp,\wp').
\]

\begin{center}
\begin{tikzpicture}[>=Latex, every node/.style={font=\small}]
  \node[visualnode, minimum width=2.6cm] (LO) at (0,1.7) {$L(O)$\\ \(1\)};
  \node[visualnode, minimum width=3.1cm] (L2) at (0,0.55) {$L(2O)$\\ \(1,\wp\)};
  \node[visualnode, minimum width=3.7cm] (L3) at (0,-0.75) {$L(3O)$\\ \(1,\wp,\wp'\)};
  \draw[->, line width=0.5pt] (LO) -- (L2);
  \draw[->, line width=0.5pt] (L2) -- (L3);
  \node[align=left] at (4.2,0.55)
    {pole orders at \(O\):\\[2pt]
     \(1:\,0\)\\
     \(\wp:\,2\)\\
     \(\wp':\,3\)};
  \draw[->, line width=0.55pt] (1.95,-0.75) -- node[above] {\(\Phi\)} (3.2,-0.75);
  \node[align=center] at (5.1,-0.75)
    {\([\wp:\wp':1]\)\\uses a basis of \(L(3O)\)};
\end{tikzpicture}
\end{center}

\section{Chapter 8}

\textbf{8.1.} Choose lifts \(p,q\in\C\) of \(P,Q\).  Then
\[
  \smap([P]-[Q])=p-q\pmod{\Lambda}.
\]
This is exactly the class \(P-Q\) in the quotient group \(T=\C/\Lambda\).

\begin{center}
\begin{tikzpicture}[>=Latex, every node/.style={font=\small}]
  \begin{scope}[scale=0.3]
    \drawprecisetorus
  \end{scope}
  \fill[black] (0.65,0.35) circle (1.5pt) node[above right] {$P$};
  \fill[black] (-0.55,-0.25) circle (1.5pt) node[below left] {$Q$};
  \draw[->, line width=0.5pt] (-0.55,-0.25) -- (0.65,0.35)
    node[midway, above] {$P-Q$};
  \node[align=center] at (3.9,0)
    {\([P]-[Q]\) has Abel sum \(P-Q\) in \(T\)};
\end{tikzpicture}
\end{center}

\textbf{8.2.} Choose lifts \(p,q,0\in\C\) of \(P,Q,O\).  The Abel sum is
\[
  p+q-2\cdot0\equiv p+q\pmod{\Lambda}.
\]
By Abel's theorem, \([P]+[Q]-2[O]\) is principal exactly when this class is
zero, namely when \(p+q\equiv0\pmod{\Lambda}\), or \(P+Q=O\) in \(T\).

\begin{center}
\begin{tikzpicture}[>=Latex, every node/.style={font=\small}]
  \node[smallvisual] (div) at (0,0)
    {$[P]+[Q]-2[O]$};
  \node[smallvisual] (sum) at (3.0,0)
    {Abel sum\\ \(P+Q\)};
  \node[smallvisual] (prin) at (6.0,0)
    {principal\\\(\Longleftrightarrow P+Q=O\)};
  \draw[->, line width=0.5pt] (div) -- (sum);
  \draw[->, line width=0.5pt] (sum) -- (prin);
\end{tikzpicture}
\end{center}

\textbf{8.3.} If \([P]-[O]=[Q]-[O]\), then \([P]-[Q]\) is principal.  Abel's
theorem gives \(P-Q=0\) in \(T\), hence \(P=Q\).

\begin{center}
\begin{tikzpicture}[>=Latex, every node/.style={font=\small}]
  \node[smallvisual] (classes) at (0,0)
    {$[P]-[O]=[Q]-[O]$};
  \node[smallvisual] (principal) at (3.25,0)
    {$[P]-[Q]$ principal};
  \node[smallvisual] (same) at (6.35,0)
    {$P-Q=0$\\ \(P=Q\)};
  \draw[->, line width=0.5pt] (classes) -- (principal);
  \draw[->, line width=0.5pt] (principal) -- node[above] {Abel} (same);
\end{tikzpicture}
\end{center}

\textbf{8.4.} Write \(D=\sum_i n_i[p_i]\) and \(D'=\sum_j m_j[q_j]\).  Then
\[
  \smap(D+D')=\sum_i n_i p_i+\sum_j m_j q_j
  =\smap(D)+\smap(D')
\]
in \(\C/\Lambda\).

\begin{center}
\begin{tikzpicture}[>=Latex, every node/.style={font=\small}]
  \node[smallvisual] (D) at (0,0.65) {$D$\\ \(\smap(D)\)};
  \node[smallvisual] (Dp) at (0,-0.65) {$D'$\\ \(\smap(D')\)};
  \node[smallvisual] (sum) at (3.2,0) {$D+D'$};
  \node[smallvisual] (image) at (6.2,0)
    {$\smap(D+D')=\smap(D)+\smap(D')$};
  \draw[->, line width=0.5pt] (D) -- (sum);
  \draw[->, line width=0.5pt] (Dp) -- (sum);
  \draw[->, line width=0.5pt] (sum) -- (image);
\end{tikzpicture}
\end{center}

\textbf{8.5.} If \(D\) is replaced by \(D+\divisor(f)\), then
\[
  \smap(D+\divisor(f))-\smap(D)=\smap(\divisor(f)).
\]
The divisor \(\divisor(f)\) is principal, so Abel's theorem gives
\(\smap(\divisor(f))=0\).  Therefore \(\smap\) is well-defined on divisor
classes in \(\Pic^0(T)\).  Exercise 8.4 gives additivity, so the induced map is
a group homomorphism.

\begin{center}
\begin{tikzpicture}[>=Latex, every node/.style={font=\small}]
  \node[visualnode] (Div) at (0,1.15) {$\Div^0(T)$};
  \node[visualnode] (Pic) at (0,-0.55) {$\Pic^0(T)=\Div^0(T)/\Prin(T)$};
  \node[visualnode] (T) at (5.0,-0.55) {$T=\C/\Lambda$};
  \node[font=\footnotesize] (Prin) at (0,-1.7) {$\Prin(T)$};
  \node[font=\footnotesize] (zero) at (5.0,-1.7) {$0\in T$};
  \draw[->, line width=0.55pt] (Div) -- node[left] {quotient} (Pic);
  \draw[->, line width=0.55pt] (Div) -- node[above] {\(\smap\)} (T);
  \draw[->, line width=0.55pt] (Pic) -- node[above] {induced \(\smap\)} (T);
  \draw[->, line width=0.45pt] (Prin) -- node[below] {\(\smap=0\) on principal divisors} (zero);
  \node[align=center] at (2.5,-2.45)
    {Abel's theorem kills principal divisors,\\so \(\smap\) descends.};
\end{tikzpicture}
\end{center}

\section{Chapter 9}

\textbf{9.1.} Homogeneous coordinates may be multiplied by \(z^3\):
\[
  [\wp:\wp':1]=[z+O(z^5):-2+O(z^4):z^3],
\]
which tends to \([0:1:0]\).

\begin{center}
\begin{tikzpicture}[>=Latex, every node/.style={font=\small}]
  \node[smallvisual] (aff) at (0,0)
    {$[\wp:\wp':1]$};
  \node[smallvisual] (scaled) at (3.0,0)
    {$[z+O(z^5):-2+O(z^4):z^3]$};
  \node[smallvisual] (limit) at (6.2,0)
    {$[0:1:0]=O$};
  \draw[->, line width=0.5pt] (aff) -- node[above] {multiply by \(z^3\)} (scaled);
  \draw[->, line width=0.5pt] (scaled) -- node[above] {\(z\to0\)} (limit);
\end{tikzpicture}
\end{center}

\textbf{9.2.} With \(x(z)=\wp(z)\) and \(y(z)=\wp'(z)\),
\[
  x(-z)=\wp(-z)=\wp(z)=x(z),
\]
while
\[
  y(-z)=\wp'(-z)=-\wp'(z)=-y(z).
\]
Thus \(z\mapsto -z\) fixes the \(x\)-coordinate and changes the sign of the
\(y\)-coordinate.

\begin{center}
\begin{tikzpicture}[>=Latex, every node/.style={font=\small}]
  \node[smallvisual] (z) at (0,0.65) {$z$};
  \node[smallvisual] (mz) at (0,-0.65) {$-z$};
  \node[smallvisual] (P) at (3.4,0.65) {$(x,y)=(\wp(z),\wp'(z))$};
  \node[smallvisual] (mP) at (3.4,-0.65) {$(x,-y)=(\wp(z),-\wp'(z))$};
  \draw[->, line width=0.5pt] (z) -- (P);
  \draw[->, line width=0.5pt] (mz) -- (mP);
  \draw[->, line width=0.5pt] (P) -- node[right] {reflect} (mP);
\end{tikzpicture}
\end{center}

\textbf{9.3.} For \(F(x,y)=y^2-(4x^3-g_2x-g_3)\), a singular affine point must
satisfy \(F=0\), \(F_y=2y=0\), and \(F_x=-12x^2+g_2=0\).  Thus \(y=0\) and
\(x\) is a repeated root of \(4x^3-g_2x-g_3\).  Conversely, a repeated root gives
such a singular point.

\begin{center}
\begin{tikzpicture}[>=Latex, every node/.style={font=\small}]
  \node[smallvisual] (sing) at (0,0) {singular point};
  \node[smallvisual] (partials) at (3.0,0)
    {$F=F_x=F_y=0$};
  \node[smallvisual] (root) at (6.0,0)
    {$y=0$\\ repeated cubic root};
  \draw[->, line width=0.5pt] (sing) -- (partials);
  \draw[->, line width=0.5pt] (partials) -- (root);
\end{tikzpicture}
\end{center}

\textbf{9.4.} Since \(\wp(-z)=\wp(z)\) and \(\wp'(-z)=-\wp'(z)\),
\[
  \Phi(-z)=[\wp(z):-\wp'(z):1]=\ominus\Phi(z).
\]
At \(O\), both sides are \(O=[0:1:0]\).

\begin{center}
\begin{tikzpicture}[>=Latex, every node/.style={font=\small}]
  \node[smallvisual] (torus) at (0,0) {$z\mapsto -z$\\on \(T\)};
  \node[smallvisual] (coords) at (3.1,0)
    {$[\wp:\wp':1]\mapsto[\wp:-\wp':1]$};
  \node[smallvisual] (cubic) at (6.45,0) {reflection\\\(\Phi(z)\mapsto\ominus\Phi(z)\)};
  \draw[->, line width=0.5pt] (torus) -- (coords);
  \draw[->, line width=0.5pt] (coords) -- (cubic);
\end{tikzpicture}
\end{center}

\textbf{9.5.} The torus notation is \(z\in T\).  The cubic notation is
\(\Phi(z)\in E(\C)\).  The divisor-class notation is
\[
  [\Phi(z)]-[O]\in\Pic^0(E).
\]
The first identification is the Weierstrass map \(\Phi\), and the second is the
Abel-Jacobi divisor-class map \(P\mapsto[P]-[O]\).

\begin{center}
\begin{tikzpicture}[>=Latex, every node/.style={font=\small}]
  \node[smallvisual] (T) at (0,0)
    {$z\in\C/\Lambda$};
  \node[smallvisual] (E) at (3.5,0)
    {$\Phi(z)\in E(\C)$};
  \node[smallvisual] (J) at (7.2,0)
    {$[\Phi(z)]-[O]\in\Pic^0(E)$};
  \draw[->, line width=0.55pt] (T) -- node[above] {\(\Phi=[\wp:\wp':1]\)} (E);
  \draw[->, line width=0.55pt] (E) --
    node[above] {$P\mapsto[P]-[O]$} (J);
  \node[align=center] at (3.6,-0.9)
    {one group element written in three languages};
\end{tikzpicture}
\end{center}

\textbf{9.6.} If \(\Phi(z)=\Phi(w)\) in the affine chart, then equality of
projective coordinates gives
\[
  \wp(z)=\wp(w),\qquad \wp'(z)=\wp'(w).
\]
For a generic \(x\)-value, the equation \(\wp(t)=\wp(z)\) has the two solutions
\[
  t\equiv z,\qquad t\equiv -z \pmod{\Lambda}.
\]
If \(w\equiv -z\), then oddness gives
\[
  \wp'(w)=\wp'(-z)=-\wp'(z).
\]
Together with \(\wp'(w)=\wp'(z)\), this forces \(\wp'(z)=0\).  Thus away from
the branch points where \(\wp'(z)=0\), the second coordinate removes the
generic \(z\) versus \(-z\) ambiguity and gives \(w\equiv z\).  The exceptional
half-period cases, where \(z\equiv -z\) modulo \(\Lambda\), and the point
\(O\) at infinity are handled by the full Weierstrass embedding theorem.  This
calculation is an injectivity check in the generic case, not a complete proof of
the theorem.

\begin{center}
\begin{tikzpicture}[>=Latex, every node/.style={font=\small}]
  \node[smallvisual] (samex) at (0,0.75)
    {$\wp(w)=\wp(z)$\\ \(w\equiv z\) or \(-z\)};
  \node[smallvisual] (samey) at (3.7,0.75)
    {$\wp'(w)=\wp'(z)$};
  \node[smallvisual] (odd) at (3.7,-0.75)
    {if \(w\equiv -z\),\\ \(\wp'(w)=-\wp'(z)\)};
  \node[smallvisual] (concl) at (7.1,0)
    {generic case:\\ \(w\equiv z\)};
  \draw[->, line width=0.5pt] (samex) -- (samey);
  \draw[->, line width=0.5pt] (samey) -- (concl);
  \draw[->, line width=0.5pt] (samex) -- (odd);
  \draw[->, line width=0.5pt] (odd) -- (concl);
  \node[align=center] at (3.6,-1.75)
    {exceptional half-periods: \(\wp'(z)=0\)};
\end{tikzpicture}
\end{center}

\section{Chapter 10}

\textbf{10.1.} In affine coordinates, write \(P=(x_0,y_0)\) and
\(Q=(x_0+h,y(x_0+h))\).  When the tangent is not vertical, the chord slope is
\[
  \frac{y(x_0+h)-y(x_0)}{h}\longrightarrow y'(x_0).
\]
For an implicit equation \(F(x,y)=0\), this derivative is
\[
  y'(x_0)=-\frac{F_x(P)}{F_y(P)}
\]
when \(F_y(P)\ne0\).  The limiting line is the tangent line.  Algebraically, its
intersection with the cubic has a double zero at \(P\), so \(P\) is counted with
intersection multiplicity \(2\).

\begin{center}
\begin{tikzpicture}[>=Latex, every node/.style={font=\small}]
  \draw[line width=0.6pt] (-1.6,-0.55) .. controls (-0.6,0.9) and (0.8,0.9) .. (1.7,-0.55);
  \fill[black] (-0.25,0.45) circle (1.5pt) node[above left] {$P$};
  \fill[black] (0.45,0.43) circle (1.5pt) node[above right] {$Q$};
  \draw[line width=0.5pt] (-1.1,0.5) -- (1.2,0.36);
  \draw[->, line width=0.5pt] (0.45,0.2) -- (-0.05,0.35) node[midway, below] {$Q\to P$};
  \draw[line width=0.55pt] (-0.9,0.62) -- (1.0,0.18);
  \node[align=center] at (4.0,0.15)
    {the chord limit is a tangent\\and counts \(P\) twice};
\end{tikzpicture}
\end{center}

\textbf{10.2.} Substitute \(y=mx+b\) into
\[
  y^2=4x^3-g_2x-g_3.
\]
The resulting cubic in \(x\) has roots \(x_P,x_Q,x_R\), counted with
multiplicity in the tangent case.  The coefficient of \(x^2\) gives
\[
  x_P+x_Q+x_R=\frac{m^2}{4}.
\]
The chord-tangent sum is the reflection of \(R\), so its \(x\)-coordinate is
\[
  x_{P\boxplus Q}=x_R=\frac{m^2}{4}-x_P-x_Q.
\]

\begin{center}
\begin{tikzpicture}[>=Latex, every node/.style={font=\small}]
  \node[smallvisual] (line) at (0,0) {$y=mx+b$};
  \node[smallvisual] (cubic) at (3.1,0)
    {cubic in \(x\)\\roots \(x_P,x_Q,x_R\)};
  \node[smallvisual] (vieta) at (6.4,0)
    {$x_P+x_Q+x_R=m^2/4$};
  \draw[->, line width=0.5pt] (line) -- (cubic);
  \draw[->, line width=0.5pt] (cubic) -- (vieta);
  \node[font=\footnotesize] at (3.2,-0.75) {substitution, then Vieta's formula};
\end{tikzpicture}
\end{center}

\textbf{10.3.} If \(P=(x,y)\), the vertical line \(X=xZ\) gives the affine
intersections \((x,y)\) and \((x,-y)\), since the equation contains \(y^2\).
In projective coordinates the same vertical line also meets the cubic at
\(O=[0:1:0]\).  Thus a vertical line meets \(P\), \(\ominus P\), and \(O\), and
the third intersection construction gives
\[
  P\boxplus(\ominus P)=O.
\]

\begin{center}
\begin{tikzpicture}[>=Latex, every node/.style={font=\small}]
  \begin{scope}[scale=0.55]
    \drawprecisecurve
    \draw[dashed, line width=0.5pt] (0.45,-2.2) -- (0.45,2.2);
    \fill[black] (0.45,1.35) circle (1.7pt) node[above right] {$P$};
    \fill[black] (0.45,-1.35) circle (1.7pt) node[below right] {$\ominus P$};
    \node at (0,-3.2) {$O$};
  \end{scope}
  \node[align=center] at (4.2,0)
    {a vertical line gives the inverse relation\\\(P\boxplus(\ominus P)=O\)};
\end{tikzpicture}
\end{center}

\textbf{10.4.} Let \(f=\ell/h=\ell/Z\) on \(E\).  The zeros of \(f\) are the
zeros of \(\ell|_E\), namely the three intersections of the line \(\ell=0\)
with the cubic:
\[
  [P]+[Q]+[R],
\]
counted with multiplicity.  The poles of \(f\) are the zeros of \(Z|_E\).  On
the Weierstrass cubic,
\[
  Z=0 \quad\Longrightarrow\quad 0=4X^3,
\]
so the line at infinity meets \(E\) only at \(O=[0:1:0]\), with total
intersection multiplicity \(3\).  Hence the pole divisor is \(3[O]\), and
\[
  \divisor(\ell/Z)=[P]+[Q]+[R]-3[O].
\]
Its degree is \(1+1+1-3=0\).

\begin{center}
\begin{tikzpicture}[>=Latex, every node/.style={font=\small}]
  \node[smallvisual] (line) at (0,0.75)
    {\(\ell|_E=0\)\\zeros \(P,Q,R\)};
  \node[smallvisual] (div) at (3.15,0.75)
    {\(Z|_E=0\)\\pole \(3[O]\)};
  \node[smallvisual] (deg) at (3.15,-0.65)
    {\(\divisor(\ell/Z)\)\\ \([P]+[Q]+[R]-3[O]\)};
  \draw[->, line width=0.5pt] (line) -- (div);
  \draw[->, line width=0.5pt] (div) -- node[right] {subtract} (deg);
\end{tikzpicture}
\end{center}

\textbf{10.5.} Since \([P]+[Q]+[R]-3[O]\) is principal, Abel's theorem says its
Abel sum is the identity.  Choose torus preimages \(p,q,r,0\) of \(P,Q,R,O\).
The Abel sum is
\[
  p+q+r-3\cdot0\equiv0\pmod{\Lambda}.
\]
In the transported group law this is exactly
\[
  P+Q+R=O.
\]

\begin{center}
\begin{tikzpicture}[>=Latex, every node/.style={font=\small}]
  \begin{scope}[scale=0.62]
    \drawprecisecurve
    \draw[black, line width=0.65pt]
      (\chordlefta,\chordleftb) -- (\chordrighta,\chordrightb);
    \fill[black] (\chordPx,\chordPy) circle (2pt);
    \node[below left] at (\chordPx,\chordPy) {$P$};
    \fill[black] (\chordQx,\chordQy) circle (2pt);
    \node[above right] at (\chordQx,\chordQy) {$Q$};
    \fill[black] (\chordRx,\chordRy) circle (2pt);
    \node[above left] at (\chordRx,\chordRy) {$R$};
    \node at (0,-3.35) {$O=[0:1:0]$};
  \end{scope}
  \node[align=center] at (5.6,0.2)
    {\([P]+[Q]+[R]-3[O]\in\Prin(E)\)\\[4pt]
     Abel sum \(=0\)\\[4pt]
     \(P+Q+R=O\)};
  \draw[->, line width=0.55pt] (2.7,0.2) -- (4.2,0.2);
\end{tikzpicture}
\end{center}

\section{Chapter 11}

\textbf{11.1.} Abel's theorem is used after the line divisor is known to be
principal.  After choosing torus preimages \(u,v,r\) of \(P,Q,R\), the Abel sum
of the divisor is
\[
  u+v+r-3\cdot0\equiv u+v+r\pmod{\Lambda}.
\]
Abel's theorem converts
\[
  [P]+[Q]+[R]-3[O]\in\Prin(E)
\]
into the group equation \(u+v+r=0\) on the torus.

\begin{center}
\begin{tikzpicture}[>=Latex, every node/.style={font=\small}]
  \node[smallvisual] (line) at (0,0)
    {$[P]+[Q]+[R]-3[O]$\\ principal};
  \node[smallvisual] (abel) at (3.4,0) {Abel's theorem};
  \node[smallvisual] (torus) at (6.5,0) {$u+v+r=0$\\on \(\C/\Lambda\)};
  \draw[->, line width=0.5pt] (line) -- (abel);
  \draw[->, line width=0.5pt] (abel) -- (torus);
\end{tikzpicture}
\end{center}

\textbf{11.2.} A direct proof must compute the third intersection of the line
through \(P,Q\), reflect it to get \(P\boxplus Q\), then compute the third
intersection of the line through \(P\boxplus Q,R\).  The other bracketing
requires the analogous line through \(Q,R\) and then the line through
\(P,Q\boxplus R\).  Comparing the final reflected points is a nonlinear
intersection calculation.  On the torus, the corresponding check is the linear
calculation
\[
  (u+v)+w=u+(v+w).
\]
The transported proof avoids the direct intersection comparison by inheriting
associativity from \(\C/\Lambda\).

\begin{center}
\begin{tikzpicture}[>=Latex, every node/.style={font=\small}]
  \node[smallvisual] (direct) at (0,0.65) {direct cubic proof\\many line intersections};
  \node[smallvisual] (transport) at (0,-0.75) {transported proof\\add in \(\C/\Lambda\)};
  \node[smallvisual] (assoc) at (4.0,-0.75) {associativity is inherited};
  \draw[->, line width=0.5pt] (transport) -- (assoc);
  \draw[dashed, line width=0.45pt] (direct) -- node[right] {harder} (transport);
\end{tikzpicture}
\end{center}

\textbf{11.3.} For \(P=(x,y)\), the projective line \(X=xZ\) contains
\[
  P=(x,y,1),\qquad \ominus P=(x,-y,1),\qquad O=[0:1:0].
\]
Indeed, \(O\) satisfies \(X=xZ\) because both sides are \(0\), and the affine
intersections have the same \(x\)-coordinate with \(y\)-coordinates \(y\) and
\(-y\).  Thus the third intersection of the line through \(O\) and \(P\) is
\(\ominus P\).  Reflecting \(\ominus P\) gives \(P\), so
\[
  O\boxplus P=P.
\]
The case \(P=O\) gives \(O\boxplus O=O\) by the tangent construction at the
identity.

\begin{center}
\begin{tikzpicture}[>=Latex, every node/.style={font=\small}]
  \node[smallvisual] (line) at (0,0)
    {projective line\\\(X=xZ\)};
  \node[smallvisual] (third) at (3.2,0)
    {points\\\(P,\ominus P,O\)};
  \node[smallvisual] (reflect) at (6.15,0)
    {reflect\\\(O\boxplus P=P\)};
  \draw[->, line width=0.5pt] (line) -- (third);
  \draw[->, line width=0.5pt] (third) -- (reflect);
\end{tikzpicture}
\end{center}

\textbf{11.4.} Let \(R=\Phi(r)\) be the third intersection of the line through
\(P=\Phi(u)\) and \(Q=\Phi(v)\).  The line divisor is principal:
\[
  [P]+[Q]+[R]-3[O].
\]
Abel's theorem gives \(u+v+r=0\) in \(\C/\Lambda\), hence
\[
  r=-(u+v).
\]
Reflecting \(R=\Phi(r)\) gives
\[
  \ominus R=\Phi(-r)=\Phi(u+v).
\]
Thus the chord-tangent construction gives the transported torus sum.

\begin{center}
\begin{tikzpicture}[>=Latex, every node/.style={font=\small}]
  \node[smallvisual] (third) at (0,0.8) {$R=\Phi(r)$\\third intersection};
  \node[smallvisual] (abel) at (3.2,0.8) {$u+v+r=0$};
  \node[smallvisual] (neg) at (6.0,0.8) {$r=-(u+v)$};
  \node[smallvisual] (ref) at (3.2,-0.75) {$\ominus R=\Phi(-r)=\Phi(u+v)$};
  \draw[->, line width=0.5pt] (third) -- (abel);
  \draw[->, line width=0.5pt] (abel) -- (neg);
  \draw[->, line width=0.5pt] (neg) -- (ref);
\end{tikzpicture}
\end{center}

\textbf{11.5.} The reflection enters after the third intersection is identified
as \(R=\Phi(r)\) with \(r=-(u+v)\).  Reflection on the cubic corresponds to
\(z\mapsto -z\) on the torus because
\[
  \Phi(-z)=\ominus\Phi(z).
\]
Therefore the reflected point is \(\Phi(-r)=\Phi(u+v)\).

\begin{center}
\begin{tikzpicture}[>=Latex, every node/.style={font=\small}]
  \node[visualnode] (uv) at (0,1.1) {$u,v\in\C/\Lambda$};
  \node[visualnode] (sum) at (4.2,1.1) {$u+v\in\C/\Lambda$};
  \node[visualnode] (PQ) at (0,-0.8) {$P=\Phi(u),\ Q=\Phi(v)$};
  \node[visualnode] (box) at (4.2,-0.8) {$P\boxplus Q$};
  \draw[->, line width=0.55pt] (uv) -- node[above] {add on torus} (sum);
  \draw[->, line width=0.55pt] (uv) -- node[left] {\(\Phi\)} (PQ);
  \draw[->, line width=0.55pt] (sum) -- node[right] {\(\Phi\)} (box);
  \draw[->, line width=0.55pt] (PQ) -- node[below] {chord--tangent} (box);
  \node[align=center] at (2.1,-1.75)
    {The square commutes: \(P\boxplus Q=\Phi(u+v)\).};
\end{tikzpicture}
\end{center}

\section{Chapter 12}

\textbf{12.1.} Under \(g_2\mapsto a^{-4}g_2\) and
\(g_3\mapsto a^{-6}g_3\), one has
\[
  g_2^3\mapsto a^{-12}g_2^3,\qquad
  g_3^2\mapsto a^{-12}g_3^2.
\]
Therefore
\[
  \Delta=g_2^3-27g_3^2\mapsto a^{-12}\Delta.
\]
The common factor \(a^{-12}\) cancels in
\[
  j=1728\frac{g_2^3}{g_2^3-27g_3^2}.
\]

\begin{center}
\begin{tikzpicture}[>=Latex, every node/.style={font=\small}]
  \node[smallvisual] (g2) at (0,0.55) {$g_2^3\mapsto a^{-12}g_2^3$};
  \node[smallvisual] (g3) at (0,-0.55) {$g_3^2\mapsto a^{-12}g_3^2$};
  \node[smallvisual] (j) at (4.4,0) {common factor cancels\\\(j\) is unchanged};
  \draw[->, line width=0.5pt] (g2) -- (j);
  \draw[->, line width=0.5pt] (g3) -- (j);
\end{tikzpicture}
\end{center}

\textbf{12.2.} In genus one, \(H^1(E,\OO_E)\cong\C\) as a complex vector space,
and \(H^1(E,\Z)\cong\Z^2\).  Under the natural map into \(H^1(E,\OO_E)\), the
integral classes give a rank-two lattice \(\Lambda\subset\C\), generated by two
periods.  Therefore
\[
  H^1(E,\OO_E)/H^1(E,\Z)\cong\C/\Lambda.
\]

\begin{center}
\begin{tikzpicture}[>=Latex, every node/.style={font=\small}]
  \node[smallvisual] (vector) at (0,0.6) {$H^1(E,\OO_E)\cong\C$};
  \node[smallvisual] (lattice) at (0,-0.7) {$H^1(E,\Z)\leadsto\Lambda$};
  \node[smallvisual] (quot) at (4.0,0) {$H^1(E,\OO_E)/H^1(E,\Z)\cong\C/\Lambda$};
  \draw[->, line width=0.5pt] (vector) -- (quot);
  \draw[->, line width=0.5pt] (lattice) -- (quot);
\end{tikzpicture}
\end{center}

\textbf{12.3.} For \(\Lambda=\Z+i\Z\), multiplication by \(i\) sends
\[
  1\mapsto i,\qquad i\mapsto -1,
\]
so it preserves the lattice and has order \(4\).  For
\(\Lambda=\Z+\tau\Z\), \(\tau=e^{2\pi i/3}\), multiplication by \(\tau\) sends
\[
  1\mapsto \tau,\qquad \tau\mapsto \tau^2=-1-\tau,
\]
so it preserves the lattice and has order \(3\).  Together with multiplication
by \(-1\), the hexagonal lattice has rotations of order \(6\).  These rotations
suggest extra automorphisms.

\begin{center}
\begin{tikzpicture}[>=Latex, every node/.style={font=\small}]
  \node[smallvisual] (square) at (0,0)
    {$\tau=i$\\square lattice\\order \(4\) rotation};
  \node[smallvisual] (hex) at (4.0,0)
    {$\tau=e^{2\pi i/3}$\\hexagonal lattice\\order \(6\) rotation};
  \draw[->, line width=0.5pt] (square) -- node[above] {extra symmetry} (hex);
\end{tikzpicture}
\end{center}

\textbf{12.4.} The quotient re-expresses the identification
\[
  \Pic^0(E)\cong\C/\Lambda.
\]
It explains the Jacobian as a complex torus.  It does not by itself prove the
chord-tangent formula.  Two further calculations are needed: first, the line
section calculation
\[
  \divisor(L|_E)=[P]+[Q]+[R]-3[O],
\]
and second, the Abel-sum calculation \(u+v+r\equiv0\pmod{\Lambda}\), which turns
that divisor into the group equation.

\begin{center}
\begin{tikzpicture}[>=Latex, every node/.style={font=\small}]
  \node[smallvisual] (coh) at (0,0.65)
    {cohomology quotient\\\(\Pic^0(E)\cong\C/\Lambda\)};
  \node[smallvisual] (line) at (4.2,0.65)
    {line divisor\\on the cubic};
  \node[smallvisual] (abel) at (4.2,-0.85)
    {Abel's theorem};
  \node[smallvisual] (law) at (0,-0.85)
    {chord--tangent law};
  \draw[->, line width=0.5pt] (coh) -- (line);
  \draw[->, line width=0.5pt] (line) -- (abel);
  \draw[->, line width=0.5pt] (abel) -- (law);
\end{tikzpicture}
\end{center}

\section{Chapter 13}

\textbf{13.1.} Over \(\F_5\), compute \(x^3+2x+1\):
\[
  1,4,3,4,3
\]
for \(x=0,1,2,3,4\).  The quadratic residues modulo \(5\) are \(0,1,4\).
Thus \(x=0\) gives two points, \(x=1\) gives two points, \(x=2\) gives none,
\(x=3\) gives two points, and \(x=4\) gives none.  Adding the point at infinity
gives \(\#E(\F_5)=7\).

\begin{center}
\begin{tikzpicture}[>=Latex, every node/.style={font=\small}]
  \foreach \x/\val/\pts in {0/1/2,1/4/2,2/3/0,3/4/2,4/3/0}{
    \node[smallvisual, minimum width=1.15cm] at (1.25*\x,0)
      {$x=\x$\\ \(f(x)=\val\)\\ \(\pts\) pts};
  }
  \node[smallvisual] at (6.9,0) {$+O$\\ total \(7\)};
  \draw[->, line width=0.5pt] (5.35,0) -- (6.1,0);
\end{tikzpicture}
\end{center}

\textbf{13.2.} The chord-tangent formulas use addition, subtraction,
multiplication, and division by a nonzero denominator.  For example, the slope is
\[
  m=\frac{y_Q-y_P}{x_Q-x_P}
\]
for a nonvertical chord, and the tangent slope is obtained by a similar quotient
involving the derivative.  Therefore the formulas make sense over any field
where the curve is smooth.  The exceptional denominator cases are vertical
lines, the point at infinity, coincident points requiring the tangent formula,
and tangent denominators equal to zero.

\begin{center}
\begin{tikzpicture}[>=Latex, every node/.style={font=\small}]
  \node[smallvisual] (smooth) at (0,0)
    {smooth cubic over \(k\)};
  \node[smallvisual] (ops) at (3.0,0)
    {field operations\\in \(k\)};
  \node[smallvisual] (law) at (6.0,0)
    {chord--tangent\\group law};
  \draw[->, line width=0.5pt] (smooth) -- (ops);
  \draw[->, line width=0.5pt] (ops) -- (law);
\end{tikzpicture}
\end{center}

\textbf{13.3.} The Hasse bound says that \(\#E(\F_p)\) lies within
\(2\sqrt p\) of \(p+1\).  Geometrically, the number of finite-field points is
close to the projective-line count \(p+1\), with a square-root-size error.

\begin{center}
\begin{tikzpicture}[>=Latex, every node/.style={font=\small}]
  \draw[->, black!55, line width=0.45pt] (-3.1,0) -- (3.3,0);
  \draw[line width=0.6pt] (-2.0,0.12) -- (2.0,0.12);
  \fill[black] (0,0.12) circle (1.5pt) node[above] {$p+1$};
  \node[below] at (-2.0,0) {$p+1-2\sqrt p$};
  \node[below] at (2.0,0) {$p+1+2\sqrt p$};
  \node[align=center] at (0,-1.05)
    {\(\#E(\F_p)\) lies in this interval};
\end{tikzpicture}
\end{center}

\textbf{13.4.} The requested separation is:
\[
\begin{array}{c|c}
\text{algebraic over a general field} & \text{analytic over }\C\\ \hline
\text{smooth cubic equation} & \text{lattice }\Lambda\subset\C\\
O=[0:1:0] & \C/\Lambda\\
\text{divisors and divisor classes} & \wp,\wp'\\
\text{rational chord--tangent formulas} & \text{period integrals}\\
E(k) \text{ for a field } k & \text{complex Lie group structure}
\end{array}
\]
The analytic uniformization
\[
  E(\C)\cong\C/\Lambda
\]
uses complex analysis, lattices in \(\C\), and meromorphic functions such as
\(\wp\); that part is special to the complex setting.

\begin{center}
\begin{tikzpicture}[>=Latex, every node/.style={font=\small}]
  \node[smallvisual] (alg) at (0,0.65)
    {algebraic cubic\\works over many fields};
  \node[smallvisual] (ct) at (4.0,0.65)
    {chord--tangent formulas};
  \node[smallvisual] (complex) at (0,-0.85)
    {analytic uniformization\\needs \(\C\), \(\Lambda\), \(\wp\)};
  \node[smallvisual] (torus) at (4.0,-0.85)
    {\(E(\C)\cong\C/\Lambda\)};
  \draw[->, line width=0.5pt] (alg) -- (ct);
  \draw[->, line width=0.5pt] (complex) -- (torus);
\end{tikzpicture}
\end{center}

\section{Chapter 14}

\textbf{14.1.} A line through \(P,Q,R\) makes
\[
  [P]+[Q]+[R]-3[O]
\]
principal.  Abel's theorem gives the group equation
\[
  P+Q+R=O.
\]
Solving for the group sum gives \(P+Q=-R=\ominus R\).  The chord--tangent rule
therefore defines \(P\boxplus Q=\ominus R\); the reflection step is forced by
the group equation.

\begin{center}
\begin{tikzpicture}[>=Latex, every node/.style={font=\small}]
  \node[smallvisual] (line) at (0,0.75)
    {line through \(P,Q,R\)};
  \node[smallvisual] (abel) at (3.15,0.75)
    {$P+Q+R=O$};
  \node[smallvisual] (reflect) at (3.15,-0.7)
    {\(P\boxplus Q=-R\)\\reflection is forced};
  \draw[->, line width=0.5pt] (line) -- (abel);
  \draw[->, line width=0.5pt] (abel) -- node[right] {reflect} (reflect);
  \node[font=\footnotesize] at (1.55,0.25) {Abel's theorem};
\end{tikzpicture}
\end{center}

\textbf{14.2.} One possible table is:
\[
\begin{array}{c|c|c}
\text{analytic} & \text{divisor-theoretic} & \text{algebraic}\\ \hline
\C/\Lambda,\wp & \Pic^0(E),\text{ divisors} & E(\C),\text{ lines}\\
\Phi=[\wp:\wp':1] & \text{Abel--Jacobi} & P\mapsto[P]-[O]
\end{array}
\]
The first map sends the torus to the cubic, and the second sends a cubic point
to its degree-zero divisor class.

\begin{center}
\begin{tikzpicture}[>=Latex, every node/.style={font=\small}]
  \node[smallvisual] (T) at (0,1.2) {analytic\\\(\C/\Lambda,\wp\)};
  \node[smallvisual] (J) at (-2.6,-0.8) {geometric\\\(\Pic^0(E)\)};
  \node[smallvisual] (E) at (2.6,-0.8) {algebraic\\\(E(\C)\), lines};
  \draw[->, line width=0.5pt] (T) -- node[above right] {\(\Phi\)} (E);
  \draw[->, line width=0.5pt] (E) -- node[below] {$P\mapsto[P]-[O]$} (J);
  \draw[->, line width=0.5pt] (J) -- node[above left] {Abel--Jacobi} (T);
\end{tikzpicture}
\end{center}

\textbf{14.3.} A compact chapter-to-method table is:
\[
\begin{array}{c|c|c}
\text{chapters} & \text{viewpoint} & \text{calculation type}\\ \hline
1\text{--}2 & \text{complex-analysis background} & \text{residues, integrals, genus}\\
3\text{--}5 & \text{analytic torus} & \text{lattice algebra, pole orders}\\
6\text{--}8 & \text{divisors and Jacobian} & \text{divisor coefficients, Abel sums}\\
9\text{--}11 & \text{plane cubic} & \text{Vieta and line intersections}\\
12 & \text{cohomology and moduli} & \text{quotients and scaling weights}\\
13 & \text{arithmetic} & \text{finite-field point counts}
\end{array}
\]

\begin{center}
\begin{tikzpicture}[>=Latex, every node/.style={font=\small}]
  \node[smallvisual] (a) at (0,0) {Ch. 3--5\\analytic};
  \node[smallvisual] (b) at (2.4,0) {Ch. 6--8\\divisors};
  \node[smallvisual] (c) at (4.8,0) {Ch. 9--11\\cubic};
  \node[smallvisual] (d) at (7.2,0) {Ch. 12--13\\extensions};
  \draw[->, line width=0.5pt] (a) -- (b);
  \draw[->, line width=0.5pt] (b) -- (c);
  \draw[->, line width=0.5pt] (c) -- (d);
\end{tikzpicture}
\end{center}

\textbf{14.4.} A compact map-check table is:

\begin{center}
\scriptsize
\setlength{\tabcolsep}{2pt}
\begin{tabular}{|L{1.15cm}|L{2.05cm}|L{1.35cm}|L{1.35cm}|L{2.35cm}|L{2.25cm}|L{2.45cm}|}
\hline
map & formula & source & target & group-law compatibility &
holomorphicity input & inverse or bijectivity input\\
\hline
\(\Phi\) &
\(z\mapsto[\wp(z):\wp'(z):1]\) &
\(\C/\Lambda\) &
\(E(\C)\) &
direct transport calculation:
\(\Phi(u+v)=\Phi(u)\boxplus\Phi(v)\) after Ex. 11.4 &
calculation with meromorphic coordinates and the projective limit at \(O\);
global extension is theorem input &
Weierstrass embedding theorem; Ex. 9.6 gives only the generic injectivity
calculation\\
\hline
\(AJ\) &
\(P\mapsto[P]-[O]\) &
\(E(\C)\) &
\(\Pic^0(E)\) &
Abel-sum calculation in Ex. 14.5 plus Abel's theorem &
accepted Jacobian/Abel-Jacobi theorem input &
Abel's theorem and genus-one representative statement\\
\hline
\end{tabular}
\end{center}

Thus the final isomorphism as complex Lie groups uses direct calculations for
the formulas, Abel sums, reflection step, and transported addition, and uses
named theorems for the global biholomorphic and Jacobian-identification inputs.

\begin{center}
\begin{tikzpicture}[>=Latex, every node/.style={font=\small}]
  \node[smallvisual] (T) at (0,0) {$\C/\Lambda$};
  \node[smallvisual] (E) at (3.0,0) {$E(\C)$};
  \node[smallvisual] (P) at (6.0,0) {$\Pic^0(E)$};
  \draw[->, line width=0.5pt] (T) -- node[above] {\(\Phi\)} (E);
  \draw[->, line width=0.5pt] (E) -- node[above] {$P\mapsto[P]-[O]$} (P);
  \node[align=center] at (3.0,-0.95)
    {both maps preserve complex structure and addition};
\end{tikzpicture}
\end{center}

\textbf{14.5.} A point \(z\in\C/\Lambda\) is written on the cubic as
\[
  \Phi(z)=[\wp(z):\wp'(z):1],
\]
and in the Jacobian as
\[
  [\Phi(z)]-[O].
\]
Adding \(z_1,z_2\) on the torus gives \(z_1+z_2\).  On the cubic this is
\[
  \Phi(z_1)\oplus\Phi(z_2)=\Phi(z_1+z_2),
\]
which equals the chord-tangent construction.  In the Jacobian it is addition of
divisor classes:
\[
  ([\Phi(z_1)]-[O])+([\Phi(z_2)]-[O]).
\]
If \(S=\Phi(z_1+z_2)\), then Abel's theorem gives
\[
  [\Phi(z_1)]+[\Phi(z_2)]-[S]-[O]\in\Prin(E)
\]
because its Abel sum is
\[
  z_1+z_2-(z_1+z_2)-0\equiv0\pmod{\Lambda}.
\]
Therefore the class of this divisor is zero in \(\Pic^0(E)\), so
\[
  [\Phi(z_1)]+[\Phi(z_2)]-2[O]=[S]-[O].
\]
Equivalently,
\[
  ([P_1]-[O])+([P_2]-[O])=[S]-[O]
\]
for \(P_i=\Phi(z_i)\) and \(S=\Phi(z_1+z_2)\).

\begin{center}
\begin{tikzpicture}[>=Latex, every node/.style={font=\small}]
  \node[visualnode] (z) at (0,0.95) {$z\in\C/\Lambda$};
  \node[visualnode] (P) at (4.6,0.95) {$\Phi(z)\in E(\C)$};
  \node[visualnode] (D) at (2.3,-0.85) {$[\Phi(z)]-[O]\in\Pic^0(E)$};
  \draw[->, line width=0.55pt] (z) -- node[above] {\(\Phi=[\wp:\wp':1]\)} (P);
  \draw[->, line width=0.55pt] (P) -- node[right] {\(\;P\mapsto[P]-[O]\)} (D);
  \draw[->, line width=0.55pt] (D) -- node[left] {Abel--Jacobi} (z);
  \node[align=center] at (2.3,-1.8)
    {one group element, three notations};
\end{tikzpicture}
\end{center}
